% =========================================================
% Computational Linear Algebra --- Vectors and Linear Combinations
% (Strang \S{}1.1 / Johnston \S{}1.1)
% =========================================================

\section{Vectors and Linear Combinations}

% ---------------------------------------------------------
\subsection*{Vectors}

\begin{propositionbox}{Definition (Vector in $\mathbb{R}^n$)}
A \emph{vector} in $\mathbb{R}^n$ is an ordered list of $n$ real numbers,
written $\mathbf{v} = (v_1, v_2, \ldots, v_n)$.
The number $n$ is the \emph{dimension} of the vector.

\medskip
\textbf{Vector addition:}
\[
  \mathbf{v} + \mathbf{w} = (v_1 + w_1,\; v_2 + w_2,\; \ldots,\; v_n + w_n).
\]
\textbf{Scalar multiplication:}
\[
  c\mathbf{v} = (cv_1,\; cv_2,\; \ldots,\; cv_n), \quad c \in \mathbb{R}.
\]
\end{propositionbox}

\begin{remark*}[Geometric meaning]
In $\mathbb{R}^2$ and $\mathbb{R}^3$, a vector is an arrow indicating displacement.
Its entries specify length and direction, but \emph{not} location ---
the same vector can be placed anywhere in space.
Addition corresponds to following one displacement and then another (head-to-tail rule).
Scalar multiplication stretches or shrinks the arrow; negative scalars reverse its direction.
\end{remark*}

\begin{theorembox}{Theorem (Properties of Vector Operations)}
For all $\mathbf{v}, \mathbf{w}, \mathbf{x} \in \mathbb{R}^n$ and $c, d \in \mathbb{R}$:
\begin{enumerate}
  \item $\mathbf{v} + \mathbf{w} = \mathbf{w} + \mathbf{v}$ \hfill (commutativity)
  \item $(\mathbf{v} + \mathbf{w}) + \mathbf{x} = \mathbf{v} + (\mathbf{w} + \mathbf{x})$ \hfill (associativity)
  \item $c(\mathbf{v} + \mathbf{w}) = c\mathbf{v} + c\mathbf{w}$ \hfill (distributivity over vector addition)
  \item $(c + d)\mathbf{v} = c\mathbf{v} + d\mathbf{v}$ \hfill (distributivity over scalar addition)
  \item $c(d\mathbf{v}) = (cd)\mathbf{v}$ \hfill (compatibility)
\end{enumerate}
\end{theorembox}

% ---------------------------------------------------------
\subsection*{Linear Combinations}

\begin{propositionbox}{Definition (Linear Combination)}
A \emph{linear combination} of vectors $\mathbf{v}_1, \mathbf{v}_2, \ldots, \mathbf{v}_k \in \mathbb{R}^n$
is any vector of the form
\[
  c_1\mathbf{v}_1 + c_2\mathbf{v}_2 + \cdots + c_k\mathbf{v}_k,
\]
where $c_1, c_2, \ldots, c_k \in \mathbb{R}$ are called the \emph{coefficients}.

A \emph{convex combination} is a linear combination in which $c_i \geq 0$
for all $i$ and $\sum_{i=1}^k c_i = 1$.
\end{propositionbox}

\begin{remark*}[Why convex combinations matter for NURBS]
Every point on a B\'{e}zier or B-spline curve is a convex combination of control points.
The Bernstein basis functions $B_{i,n}(u)$ are non-negative and sum to 1,
so $C(u) = \sum_{i=0}^n B_{i,n}(u) P_i$ is a convex combination of the control
points $\{P_i\}$ at every parameter value $u$.
This is the \emph{convex hull property}: the curve lies entirely within the convex hull
of its control polygon.
\end{remark*}

\begin{propositionbox}{Definition (Standard Basis Vectors)}
For $j = 1, 2, \ldots, n$, the \emph{$j$-th standard basis vector} $\mathbf{e}_j \in \mathbb{R}^n$
has a $1$ in position $j$ and $0$s everywhere else.
Every vector $\mathbf{v} = (v_1, \ldots, v_n) \in \mathbb{R}^n$ can be written as
\[
  \mathbf{v} = v_1 \mathbf{e}_1 + v_2 \mathbf{e}_2 + \cdots + v_n \mathbf{e}_n.
\]
\end{propositionbox}
